\chapter{Разговор через воздух}
% полная версия: chapters/05_serocomputer.tex

\pencilfig{5}{\pencilart{5}}{Слева телефон: провод от одного к одному. Справа радио: одно сообщение сразу всем вокруг.}

До сих пор нейроны говорили друг с другом по телефону: провод к проводу, адресно. Теперь включим первую радиостанцию --- \nt{СЕРО}-нейроны, выбрасывающие серотонин. Их всего несколько сотен тысяч на весь мозг, но их провода расходятся по огромным его областям.

\section*{Четыре отличия}

\nt{СЕРО}-нейрон непохож на \nt{ГЛУТ}. Во-первых, знак его сигнала выбирает получатель: у серотонина есть «замки» и того, и другого знака. Во-вторых, он работает сам: в бодрствовании ровно щёлкает несколько раз в секунду, если получает постоянную фоновую подпитку, --- это метроном. В-третьих, он медленный: его сигнал действует не за миллисекунды, а за сотни миллисекунд. В-четвёртых, он часто выбрасывает вещество не в узкую щель, а в пространство вокруг, где оно расплывается, как капля чернил в воде. Получатель чувствует концентрацию и не знает, от кого пришла молекула.

\section*{Логика, которой нет}

Метроном с тормозящим «замком» --- это готовое «не»: пока серотонина нет, нейрон работает; появился --- замолчал. Один нейрон вместо целой схемы. Кажется, серотониновая сеть логически сильнее \nt{ГЛУТ}-сети.

Но есть загвоздка. Выброшенное в пространство вещество смешивается. Два сигнала остаются разными, только если их выпустили в местах, разнесённых дальше, чем вещество успевает расплыться, --- а это десятки микрометров. У каждого же серотонинового нейрона огромное количество мест выброса, раскиданных по большим областям, и «провода» разных нейронов перекрываются. Отдельных битов по такой сети не передашь. Она может передавать лишь несколько общих величин --- как радио, по которому все слушают одну и ту же станцию.

\pencilfig{123}{%
% two release sites: far apart the clouds stay separate (two signals), close together they merge (one level)
\foreach \x/\y in {0.9/1.55, 4.1/1.55, 1.9/0, 2.9/0}{
  \foreach \r/\o in {0.95/0.07, 0.7/0.09, 0.45/0.12}{\fill[pcViolet, opacity=\o] (\x,\y) circle (\r);}
  \draw[dotpen=pcViolet, fill=pcViolet] (\x,\y) circle (0.07);}
\draw[arr] (1.95,1.55) -- (3.05,1.55); \draw[arr] (3.05,1.55) -- (1.95,1.55);
\node[lbl, anchor=south, align=center] at (2.5,1.6) {дальше, чем\\расплывается};
\node[lbl, anchor=west] at (5.25,1.55) {два сигнала};
\node[lbl, anchor=west] at (5.25,0) {одно общее облако};
}{Серотонин выбрасывается в пространство и расплывается облаком. Два источника дают два разных сигнала, только если они дальше друг от друга, чем вещество успевает расплыться. Если ближе, облака сливаются в один общий уровень.}

\section*{Что она делает на самом деле}

Для одной общей величины у этой системы всё есть. Концентрация серотонина сглаживает отдельные щелчки в плавный уровень --- так простейший преобразователь превращает поток импульсов в напряжение. Этот уровень держится секунды и сам тает: не бит, который надо стирать, а медленно исчезающий след.

Ещё у серотониновых нейронов есть «замки» на самих себе: выброшенный серотонин притормаживает собственный источник. Инженер узнает в этом усилитель с отрицательной обратной связью --- термостат, держащий уровень. Правда, опыты показывают, что это самоторможение устроено адресно и вызывает лишь короткие паузы, так что роль термостата пока скорее гипотеза, чем факт.

\section*{Сколько стоит ждать}

Какую общую величину стоит держать медленной и одинаковой для всей сети? Хороший кандидат --- \emph{терпение}. Представьте выбор: маленькая награда сейчас или большая, но через несколько секунд. Сколько вы готовы ждать, зависит от того, как быстро для вас «обесценивается» будущее. По одной из гипотез, эту ручку и крутит серотонин. Есть опыты в её пользу: если искусственно включить серотониновые нейроны, мышь дольше ждёт отложенную награду, прежде чем сдаться.

Эта глава вводит три устройства, которыми будут пользоваться все радиостанции мозга: нейрон-метроном, «провод», который на самом деле объём, и медленный уровень вместо отдельных щелчков.

\fullversion{5}
